\subsection{Radiative decay of \texorpdfstring{$K^+ \rightarrow e^+ \nu \gamma$}{K+ -> e+ nu g} with TREK/E36 at J-PARC}


The TREK collaboration (Time Reversal Experiment with Kaons)
was formed in 2005 to pursue a program to search for T-violation in stopped
$K^+$ decays~\cite{e06_proposal}.
When by 2010 it became clear that the
required kaon beam intensity would not be achievable in a timely manner, the
TREK collaboration still pursued a program on sensitive
beyond-the-Standard-Model physics, with experiment E36, which was proposed in
2010 using only 50 kW beam power~\cite{e36_proposal,e36_proposal-addend,e36_proposal-addend2}.

The focus of E36 was on testing lepton universality in the
SM, by a precision measurement of the ratio of two-body decay widths,
$R_K = \Gamma(K_{e2})/\Gamma(K_{\mu2})$.
The experiment was approved in 2013 and constructed at the J-PARC K1.1BR
beamline in 2014, reusing the toroidal spectrometer setup from KEK-PS
experiment E246~\cite{e246_NIM} and the CsI(Tl) barrel~\cite{ito_csi}.
A new scintillating fiber target,
a spiral fiber tracker~\cite{sft}, and
redundant PID systems (an aerogel counter array~\cite{ac},
a time-of-flight system, and
a leadglass shower calorimeter in each sector~\cite{pgc}) were implemented.
Data taking was limited to less than three months in 2015. First E36 results
have been published on radiative kaon
decays~\cite{e36_Ke2g_CsI,e36_Ke2g_GSC},
with the result for $R_K$ being expected to become available in 2023.
The E36 setup was also used for an in-situ search for dark photons, or more
generally for light neutral bosons decaying into $e^+ e^-$ pairs over a broad
mass range of 10--200\;MeV/c$^2$.
The analysis, which is still ongoing, combines a charged particle
($\mu^+$, $\pi^+$, or $e^+$) in one toroidal sector with two clusters from
$e^+ e^-$ pair detection in the segmented CsI(Tl) calorimeter.
Preliminary results from the light boson search have been
reported~\cite{e36_dp}, and the
final sensitivity is expected to be reached in 2023.

The study of radiative kaon decay allowed to extract the branching ratio for
the structure-dependent (SD) process resulting in a hard photon radiated nearly
in back-to-back kinematics with the positron. The process is accompanied by the
internal bremsstrahlung (IB), predominantly of soft photons nearly in the
direction of the positron. In the overlapping regions both processes appear,
and the IB process is calculated and simulated as a background, and
conventionally included in the two-body branching ratio for $K_{e2}$.
The latest analysis also considers internal bremsstrahlung in the
SD process $K_{e2\gamma}^\mathrm{SD}$.
The implementation of the IB radiative process in the MC simulation
follows the scheme introduced by Gatti~\cite{gatti2006}, allowing for analytic
expressions of the differential branching ratio distributions to be used in an
event generator. The experiment data were
analyzed in terms of the lepton momentum in a narrow range. The radiative
correction affecting the $K_{e2}$ and $K_{e2\gamma}$
  branching ratios were mostly due to the narrow acceptance chosen for the
  lepton momentum. In the ratio
$R_\gamma = \Gamma(K_{e2\gamma}^\mathrm{SD})/\Gamma(K_{e2})$,
the IB corrections partially cancel in the acceptance ratio.
The first publication~\cite{e36_Ke2g_CsI}
accounted for the IB acceptance effect only in the
monochromatic two-body decay $K_{e2}$. In the most recent
result~\cite{e36_Ke2g_GSC}, the IB scheme was also
implemented to correct the radiative
yield in the three-body decay $K_{e2\gamma}$.
The obtained ratio is consistent with
a recent lattice-QCD calculation~\cite{latticeQCD}, a gauged nonlocal chiral
quark model (NL$\chi$QM)~\cite{NLchiQM},
however at tension with ChPT at order $\mathcal{O}(p^4)$~\cite{bijnens1993,daphne1995}
and $\mathcal{O}(p^6)$~\cite{chen2008}, and a previous
measurement by KLOE~\cite{KLOE} at the 4-$\sigma$ level.
The latter data have been
used by NA62~\cite{NA62} to correct the universality ratio $R_K$ for the SD
background, which should be revisited.

\subsection{Radiative decay width \texorpdfstring{$\rho^- \rightarrow \pi^- \gamma$}{rho- -> pi- g} with COMPASS at CERN}

The COMPASS collaboration (COmmon Muon and Proton Apparatus for Structure and Spectroscopy) operated a multi-purpose,
high-resolution and high-rate capable spectrometer~\cite{Abbon_2007} at the CERN SPS M2 beamline with muon and hadron beams in the energy range 100--280\;GeV. A program pursued since the very beginning some 20 years ago focuses on ultra-peripheral reactions of pion and kaon beams on atomic nuclei, in the domain of momentum transfers where the electromagnetic interaction dominates --- so-called Primakoff reactions. It is intended to continue the program with a focus on kaon-induced reaction with AMBER~\cite{Adams:2018pwt}, operating at the same site as COMPASS and as its successor experiment.
\begin{figure}[ht]
    \centering
    \includegraphics[width=\columnwidth]{Mesons/fig/ChiralAnoRho.pdf}
    \caption{Contributions to the cross section for the reaction $\pi^-\gamma\rightarrow\pi^- \pi^0$ from the kinematic threshold, where the chiral anomaly dominates, to the region of the $\rho(770)$ resonance. Adapted from the respective figure in~\cite{Kaiser_2008}. The simplified model is used here to display the conceptual structure; for the data analysis to extract the chiral anomaly and the radiative width of the resonance, the dispersive framework~\cite{PhysRevD.86.116009} is used.}
    \label{fig:ChARho}
\end{figure}

One of the channels of interest is $\pi^-\gamma^{(*)}\rightarrow\pi^- \pi^0$, the photon from the exchange with the nucleus being quasi-real. According to the equivalent-photon formalism of Weizs\"acker and Williams, for a 190\;GeV pion beam on nickel, these photons cover a wide range of momenta of 0.1 to 2.0\;GeV in the rest frame of the incoming pion. Thus the full low-momentum range for the reaction becomes accessible and covers the kinematic region in which the $\rho(770)$ resonance dominates in the $s$-channel, $\pi^-\gamma^{(*)}\rightarrow\rho^-\rightarrow\pi^- \pi^0$. The cross section in the low-energy region, governed by the chiral anomaly up to the resonance region, is depicted in Fig.~\ref{fig:ChARho}. The strength of the resonance, as it appears in the reaction, depends on the radiative coupling $\rho\rightarrow\pi\gamma$, entering in inverted kinematics. The measurement of the cross section thus allows for a precision determination of the radiative width of the $\rho(770)$.

The reaction is affected by higher-order electromagnetic processes that have to be corrected for. A dominant contribution comes from a process with a $t$-channel exchange photon, coupling with the incoming photon to the outgoing $\pi ^0$ (thus $\pi^0\rightarrow\gamma\gamma$ in inverted kinematics). Corrections due to photon loops and real-photon emission have been calculated and found to be on the level of a few percent, thus are to be taken into account for a precision analysis. The formulae have been given in~\cite{PhysRevD.64.094009} and revised \cite{Kaiserpriv}. The full implementation in terms of an event generator including bremsstrahlung photons has been started.